\documentclass[11pt,a4paper]{article}

% ---------- Packages ----------
\usepackage[utf8]{inputenc}
\usepackage[T1]{fontenc}
\usepackage[italian]{babel}
\usepackage[a4paper,margin=2.5cm]{geometry}
\usepackage{graphicx}
\usepackage{booktabs}
\usepackage{array}
\usepackage{longtable}
\usepackage{listings}
\usepackage{xcolor}
\usepackage{amssymb}
\usepackage{hyperref}
\usepackage{tikz}
\usepackage{enumitem}
\usetikzlibrary{positioning, arrows.meta, shapes.geometric, fit}

\hypersetup{
  colorlinks=true,
  linkcolor=black!50!black,
  urlcolor=blue!50!black,
  citecolor=blue!50!black,
  pdftitle={OpenTeX -- Relazione finale},
}

\lstdefinestyle{code}{
  basicstyle=\ttfamily\small,
  breaklines=true,
  frame=single,
  backgroundcolor=\color{gray!8},
  columns=fullflexible,
  keywordstyle=\color{blue!70!black}\bfseries,
  commentstyle=\color{green!40!black},
  showstringspaces=false,
}
\lstset{style=code}

\newcommand{\Done}{\textcolor{green!50!black}{\checkmark}}

\title{\textbf{OpenTeX}\\[0.3em]\large Ambiente collaborativo LaTeX self-hosted su MongoDB\\[0.3em]}
\author{
  Cieri Manuel, Citro Gianluigi
% TODO: Team ID
}
\date{
 %28/07/2026
}

\begin{document}

\maketitle
\vspace{-2.5em}
\begin{abstract}
\vspace{-0.5em}
\noindent
OpenTeX è un ambiente collaborativo self-hosted per la scrittura di documenti LaTeX, sviluppato come progetto per il corso di Basi di Dati 2 (NoSQL). Il sistema adotta MongoDB come motore di persistenza, un backend FastAPI/Motor asincrono e un frontend React/Vite, con compilazione dei documenti tramite il motore \textit{Tectonic}. Questa relazione descrive l'installazione, l'architettura, il modello dei dati, le scelte CAP/BASE, le operazioni CRUD, la generazione di dati sintetici (ETL), gli indici e i relativi benchmark, e la pipeline di aggregazione (JOIN) realizzata sul progetto, riportando esclusivamente metriche misurate sul sistema realmente implementato.
\end{abstract}

\tableofcontents
\newpage

% =====================================================
\section{Introduzione}
% =====================================================

OpenTeX nasce per dimostrare l'uso di un database NoSQL documentale (MongoDB) in un caso d'uso realistico: un editor collaborativo di documenti LaTeX con gestione di utenti, progetti, permessi, file e log di attività. Il progetto è stato sviluppato in maniera incrementale attraverso 14 issue (vedi Sezione~\ref{sec:conclusioni}), ciascuna chiusa con una pull request dedicata e revisionata da un membro del team diverso dall'autore.

\subsection{Stack tecnologico}

\begin{table}[h!]
\centering
\begin{tabular}{@{}lll@{}}
\toprule
\textbf{Componente} & \textbf{Tecnologia} & \textbf{Ruolo} \\
\midrule
Frontend & React 19 + Vite 8 & Single-page application \\
HTTP client & Axios 1.x & Iniezione automatica del Bearer JWT \\
Editor & @uiw/react-codemirror 4.x & Editor con syntax highlighting LaTeX \\
Backend & FastAPI + Motor (async) & API REST e logica applicativa \\
Database & MongoDB 8.0 & Persistenza documentale \\
Compilazione & Tectonic & Compilazione LaTeX $\to$ PDF, imbarcato nel backend \\
Orchestrazione & Docker Compose & 3 servizi: mongo, backend, frontend \\
\bottomrule
\end{tabular}
\caption{Stack tecnologico di OpenTeX.}
\end{table}

% =====================================================
\section{Installazione}
% =====================================================

Il sistema è interamente containerizzato; l'unico prerequisito è Docker (\textgreater= 24.0) e Docker Compose (\textgreater= 2.20).

\subsection{Procedura}

\begin{enumerate}[leftmargin=2.2em]
  \item Clonare il repository e configurare le variabili d'ambiente:
\begin{lstlisting}
git clone <url-repo> && cd opentex
cp .env.example .env   # impostare SECRET_KEY, MONGO_ROOT_USER, MONGO_ROOT_PASSWORD
\end{lstlisting}
  \item Avviare tutti i servizi:
\begin{lstlisting}
docker compose up -d --build
\end{lstlisting}
  \item Applicare la validazione dello schema MongoDB (\texttt{\$jsonSchema}):
\begin{lstlisting}
docker compose exec backend python -m scripts.validation.apply_schema_validation
\end{lstlisting}
  \item (Opzionale) Caricare dati sintetici di esempio:
\begin{lstlisting}
python seed/seed.py --users 50 --projects 100 --logs 30000 --drop
\end{lstlisting}
  \item Verificare il funzionamento:
\begin{lstlisting}
curl http://localhost:8000/health
\end{lstlisting}
\end{enumerate}

\subsection{Servizi esposti}

\begin{table}[h!]
\centering
\begin{tabular}{@{}lll@{}}
\toprule
\textbf{Servizio} & \textbf{URL} & \textbf{Note} \\
\midrule
Frontend & \url{http://localhost:5173} & React + Vite, hot-reload \\
Backend API & \url{http://localhost:8000} & FastAPI + Motor \\
Swagger UI & \url{http://localhost:8000/docs} & Documentazione API interattiva \\
MongoDB & \texttt{localhost:27017} & Credenziali in \texttt{.env} \\
\bottomrule
\end{tabular}
\caption{Servizi e porte esposte da Docker Compose.}
\end{table}

Login di test (creato dal seed): \texttt{admin@opentex.org} / \texttt{password} (account amministratore di sito, con accesso a statistiche e benchmark). Tutti gli account generati dal seed condividono la stessa password; gli indirizzi di esempio sono riportati in \texttt{seed/seed\_output.txt} dopo ogni esecuzione.

% =====================================================
\section{Architettura}
% =====================================================

\subsection{Architettura a tre livelli}

\begin{figure}[h!]
\centering
\begin{tikzpicture}[
  box/.style={draw, rounded corners, minimum width=3.4cm, minimum height=1.6cm, align=center, fill=blue!6},
  arr/.style={-Stealth, thick}
]
\node[box] (fe) {Frontend\\React + Vite\\porta 5173};
\node[box, right=3.2cm of fe] (be) {Backend\\FastAPI + Motor\\porta 8000};
\node[box, right=3.2cm of be] (db) {MongoDB 8.0\\(Docker)\\porta 27017};

\draw[arr] (fe.east) -- (be.west) node[midway, above, font=\scriptsize] {HTTP/REST};
\draw[arr] (be.west) -- (fe.east) node[midway, below, font=\scriptsize] {JSON / PDF};
\draw[arr] (be.east) -- (db.west) node[midway, above, font=\scriptsize] {Motor (async)};
\draw[arr] (db.west) -- (be.east) node[midway, below, font=\scriptsize] {Documenti BSON};
\end{tikzpicture}
\caption{Architettura a tre livelli di OpenTeX. Tutti i servizi sono orchestrati via Docker Compose; il backend imbarca un binario statico di Tectonic, eliminando la necessità di un'installazione TeX separata.}
\end{figure}

\subsection{Modello dei dati: 5 collezioni}

Il database è organizzato in cinque collezioni. La scelta tra \emph{embedding} e \emph{referencing} è stata guidata da tre criteri: stabilità del dato, dimensione, e limite di 16\,MB per documento imposto da MongoDB.

\begin{table}[h!]
\centering
\begin{tabular}{@{}p{4.2cm}p{2.2cm}p{7cm}@{}}
\toprule
\textbf{Decisione} & \textbf{Scelta} & \textbf{Motivazione} \\
\midrule
Metadati di progetto (tag, abstract) & Embedded & Piccoli, stabili, letti sempre insieme al progetto \\
File, permessi, log di attività & Referenced (\texttt{ObjectId}) & Crescita illimitata; i file multimediali eccedono il limite di 16\,MB/documento \\
Profili utente & Referenced & Condivisi tra molti progetti; aggiornati indipendentemente \\
\bottomrule
\end{tabular}
\caption{Criteri di modellazione: embedding vs.\ referencing.}
\end{table}

\begin{table}[h!]
\centering
\small
\begin{tabular}{@{}p{2.2cm}p{9.8cm}@{}}
\toprule
\textbf{Collezione} & \textbf{Campi principali} \\
\midrule
\texttt{users} & \texttt{email} (univoca), \texttt{first\_name}, \texttt{last\_name}, \texttt{department}, \texttt{hashed\_password}, \texttt{is\_admin}, \texttt{created\_at}, \texttt{preferences} \\
\texttt{projects} & \texttt{title}, \texttt{abstract}, \texttt{tags[]}, \texttt{owner\_id} $\to$ \texttt{users}, \texttt{created\_at}, \texttt{updated\_at} \\
\texttt{files} & \texttt{project\_id} $\to$ \texttt{projects}, \texttt{filename}, \texttt{file\_type} (\texttt{tex|bib|image|pdf}), \texttt{content} (tex/bib) | \texttt{path} (image/pdf), \texttt{uploaded\_by} $\to$ \texttt{users} \\
\texttt{permissions} & \texttt{user\_id} $\to$ \texttt{users}, \texttt{project\_id} $\to$ \texttt{projects}, \texttt{role} (\texttt{Admin|Editor|Viewer}), \texttt{granted\_at}, \texttt{granted\_by} \\
\texttt{activity\_logs} & \texttt{user\_id} $\to$ \texttt{users}, \texttt{project\_id} $\to$ \texttt{projects}, \texttt{action}, \texttt{resource}, \texttt{resource\_id}, \texttt{timestamp} \\
\bottomrule
\end{tabular}
\caption{Schema sintetico delle 5 collezioni. La validazione \texttt{\$jsonSchema} applicata in MongoDB impone i campi obbligatori, i tipi e le regole condizionali (es.\ \texttt{content} richiesto per \texttt{tex}/\texttt{bib}, \texttt{path} per \texttt{image}/\texttt{pdf}); non viene invece verificata l'integrità referenziale a livello di database, scelta coerente con il modello AP descritto in Sezione~\ref{sec:cap}.}
\end{table}

\pagebreak

\subsection{Flusso di autenticazione}

L'autenticazione è basata su JWT firmati con algoritmo HS256. Le password sono salvate con hash \texttt{bcrypt} e non vengono mai restituite in alcuna risposta API. Il token scade dopo \texttt{JWT\_EXPIRATION\_MINUTES} (default 1440 minuti, 24\,h).

\begin{lstlisting}
POST /auth/register {email, password, first_name, last_name, department}
  -> bcrypt.hash(password) salvato in users; 201 Created

POST /auth/login {email, password}
  -> bcrypt.verify; JWT firmato con SECRET_KEY
  -> 200 {access_token, token_type, user}

GET /projects/  (Authorization: Bearer <token>)
  -> JWT decodificato, user_id iniettato come dipendenza
  -> 200 [...]
\end{lstlisting}

\subsection{Pipeline di compilazione LaTeX}

La compilazione avviene in una directory temporanea isolata, con timeout e parallelismo limitato per contenere il carico sul backend:

\begin{enumerate}[leftmargin=2.2em]
  \item Recupero di tutti i file \texttt{.tex}/\texttt{.bib} del progetto da MongoDB.
  \item Scrittura in una directory temporanea isolata (\texttt{tempfile.mkdtemp}).
  \item Acquisizione di un semaforo (massimo 2 job paralleli).
  \item Esecuzione di \texttt{tectonic --keep-logs --outdir <tmp> <entry\_point.tex>}, timeout 60\,s, \texttt{\textbackslash write18} disabilitato di default.
  \item Successo $\to$ lettura di \texttt{output.pdf}, risposta \texttt{application/pdf}.
  \item Fallimento $\to$ risposta \texttt{422} con il log di Tectonic.
  \item Timeout $\to$ risposta \texttt{504}.
  \item Sempre $\to$ \texttt{shutil.rmtree(tmp)} (pulizia).
\end{enumerate}

L'entry point è \texttt{main.tex} se presente, altrimenti il primo file \texttt{.tex} in ordine alfabetico.

Le tre componenti principali del flusso -- query al database (passo 1), scrittura su filesystem (passo 2) ed esecuzione del compilatore (passo 4) -- sono strumentate e misurate separatamente; la relativa analisi è riportata nella Sezione~\ref{sec:compile-bench}.

% =====================================================
\section{Modello CAP e proprietà BASE}
\label{sec:cap}
% =====================================================

OpenTeX gira su una singola istanza/replica MongoDB. Il sistema è progettato secondo il paradigma \textbf{AP} del teorema CAP:

\begin{itemize}[leftmargin=2.2em]
  \item \textbf{Availability} -- le richieste di lettura e scrittura hanno successo anche in presenza di guasti transitori.
  \item \textbf{Partition Tolerance} -- il sistema continua a operare in presenza di partizioni di rete.
  \item \textbf{Consistency} (rilassata) -- MongoDB garantisce solo atomicità a livello di singolo documento; non esistono transazioni multi-documento o cross-collection. Sono quindi possibili letture stale in casi limite (es. un documento \texttt{projects} aggiornato mentre viene letto il relativo \texttt{permissions}).
\end{itemize}

\begin{table}[h!]
\centering
\begin{tabular}{@{}p{3.4cm}p{8.6cm}@{}}
\toprule
\textbf{Proprietà BASE} & \textbf{Implementazione in OpenTeX} \\
\midrule
\textbf{B}asically Available & FastAPI risponde immediatamente; nessun lock globale \\
\textbf{S}oft state & Gli aggregati derivati (es.\ conteggio collaboratori) sono calcolati on-demand, non cacheati \\
\textbf{E}ventually consistent & I log di attività e le modifiche ai permessi si propagano alla lettura successiva; nessuna sincronizzazione write-ahead \\
\bottomrule
\end{tabular}
\caption{Semantica BASE adottata da OpenTeX.}
\end{table}

% =====================================================
\section{API e operazioni CRUD}
% =====================================================

Il backend espone 7 router FastAPI con operazioni CRUD complete sulle risorse principali, documentate automaticamente via Swagger (\texttt{/docs}).

\begin{table}[h!]
\centering
\scriptsize
\begin{tabular}{@{}p{1.9cm}p{1.1cm}p{5.4cm}p{4.4cm}@{}}
\toprule
\textbf{Router} & \textbf{Metodo} & \textbf{Endpoint} & \textbf{Descrizione} \\
\midrule
\texttt{auth} & POST & \texttt{/auth/register} & Registrazione utente \\
 & POST & \texttt{/auth/login} & Login, rilascio JWT \\
 & GET & \texttt{/auth/departments} & Elenco dipartimenti \\
\texttt{projects} & POST & \texttt{/projects/} & Creazione progetto \\
 & GET & \texttt{/projects/} & Elenco progetti (filtro \texttt{owner\_id}) \\
 & GET & \texttt{/projects/\{id\}} & Dettaglio progetto \\
 & PUT & \texttt{/projects/\{id\}} & Aggiornamento progetto \\
 & DELETE & \texttt{/projects/\{id\}} & Cancellazione progetto \\
 & GET & \texttt{/projects/\{id\}/files} & File collegati al progetto \\
 & POST & \texttt{/projects/\{id\}/files} & Creazione file nel progetto \\
\texttt{files} & PUT & \texttt{/files/\{id\}} & Salvataggio contenuto file \\
 & DELETE & \texttt{/files/\{id\}} & Cancellazione file \\
\texttt{permissions} & POST & \texttt{/projects/\{id\}/permissions} & Assegna/aggiorna ruolo \\
 & DELETE & \texttt{/projects/\{id\}/permissions/\{uid\}} & Revoca accesso \\
 & GET & \texttt{/projects/\{id\}/permissions} & Elenco collaboratori \\
\texttt{compile} & POST & \texttt{/projects/\{id\}/compile} & Compilazione LaTeX $\to$ PDF \\
\texttt{stats} & GET & \texttt{/stats/projects} & Statistiche aggregate (JOIN) \\
 & GET & \texttt{/stats/benchmarks} & Risultati benchmark indici \\
 & GET & \texttt{/stats/compile-benchmark} & Ripartizione tempi di compilazione \\
\texttt{users} & GET & \texttt{/users/} & Elenco utenti \\
 & GET & \texttt{/users/\{id\}} & Dettaglio utente \\
\bottomrule
\end{tabular}
\caption{Endpoint REST implementati (21 endpoint su 7 router).}
\end{table}

\newpage

\subsection{Controllo accessi (RBAC)}

\begin{table}[h!]
\centering
\begin{tabular}{@{}lcccc@{}}
\toprule
\textbf{Ruolo} & \textbf{Read} & \textbf{Edit} & \textbf{Delete} & \textbf{Gestione permessi} \\
\midrule
Admin  & \Done & \Done & \Done & \Done \\
Editor & \Done & \Done & --    & --    \\
Viewer & \Done & --    & --    & --    \\
\bottomrule
\end{tabular}
\caption{Ruoli di progetto (collezione \texttt{permissions}). Il proprietario del progetto è sempre Admin implicito ed è l'unico che può cancellarlo. Il flag \texttt{is\_admin} a livello account è indipendente da questi ruoli.}
\end{table}

\pagebreak

% =====================================================
\section{Generazione di dati sintetici}
% =====================================================

Per testare query, indici e benchmark in condizioni realistiche è stato sviluppato un generatore di dati sintetici basato su \texttt{Faker} (\texttt{seed/seed.py}), parametrico nel numero di utenti, progetti e log generati, e con integrità referenziale garantita (ogni \texttt{owner\_id} punta a un utente esistente, ogni \texttt{permission} a una coppia utente-progetto valida).

\begin{lstlisting}
pip install -r seed/requirements.txt
python seed/seed.py --users 50 --projects 100 --logs 30000 --drop
\end{lstlisting}

\subsection{Output di un'esecuzione reale}

\begin{table}[h!]
\centering
\begin{tabular}{@{}lr@{}}
\toprule
\textbf{Collezione} & \textbf{Documenti generati} \\
\midrule
\texttt{users} & 51 (incluso l'account admin) \\
\texttt{projects} & 100 \\
\texttt{files} & 337 \\
\texttt{permissions} & 251 \\
\texttt{activity\_logs} & 30\,000 \\
\midrule
\textbf{Tempo totale} & 0.3\,s \\
\bottomrule
\end{tabular}
\caption{Risultato di un'esecuzione dello script di seed (\texttt{seed/seed\_output.txt}).}
\end{table}

% =====================================================
\section{Indici e benchmark}
% =====================================================

\subsection{Indici creati}

\begin{table}[h!]
\centering
\scriptsize
\begin{tabular}{@{}p{4.4cm}p{2.6cm}p{3cm}p{2cm}@{}}
\toprule
\textbf{Indice} & \textbf{Collezione} & \textbf{Campi} & \textbf{Tipo} \\
\midrule
\texttt{projects\_text\_search} & \texttt{projects} & \texttt{title}, \texttt{abstract} & Full-text \\
\texttt{projects\_owner\_date} & \texttt{projects} & \texttt{owner\_id} + \texttt{created\_at} & Composto \\
\texttt{permissions\_project\_id} & \texttt{permissions} & \texttt{project\_id} & Singolo campo \\
\texttt{activity\_logs\_project\_id} & \texttt{activity\_logs} & \texttt{project\_id} & Singolo campo \\
\texttt{files\_project\_id} & \texttt{files} & \texttt{project\_id} & Singolo campo \\
\midrule
\texttt{email\_1} & \texttt{users} & \texttt{email} & Univoco \\
\bottomrule
\end{tabular}
\caption{Indici del sistema. I primi cinque sono creati da \texttt{scripts/indexes/create\_indexes.py} e verificati con \texttt{explain()}; l'indice univoco su \texttt{users.email} è garantito all'avvio dell'applicazione (\texttt{backend/app/main.py}) e assolve una duplice funzione: accelera la ricerca dell'utente in fase di login e impone a livello di database l'unicità dell'indirizzo email in fase di registrazione.}
\end{table}

\subsection{Risultati del benchmark}

Le query sono state misurate con e senza l'indice corrispondente, su 10 esecuzioni per query (media riportata).

\begin{table}[h!]
\centering
\small
\begin{tabular}{@{}p{4.6cm}rrr@{}}
\toprule
\textbf{Query} & \textbf{Con indice} & \textbf{Senza indice} & \textbf{Speedup} \\
\midrule
Text search (title + abstract) & 1.753\,ms & 2.003\,ms & 1.1$\times$ \\
Compound: \texttt{owner\_id} + \texttt{created\_at} & 1.492\,ms & 1.655\,ms & 1.1$\times$ \\
Permissions by \texttt{project\_id} & 1.561\,ms & 1.667\,ms & 1.1$\times$ \\
Activity logs by \texttt{project\_id} & 1.750\,ms & 19.024\,ms & \textbf{10.9$\times$} \\
Files by \texttt{project\_id} & 1.609\,ms & 1.979\,ms & 1.2$\times$ \\
\bottomrule
\end{tabular}
\caption{Tempi medi di esecuzione su 10 run (\texttt{scripts/benchmark/run\_benchmark.py}).}
\end{table}

\textbf{Interpretazione.} Lo speedup più marcato (10.9$\times$) si osserva sulla query su \texttt{activity\_logs}, l'unica collezione realmente di grandi dimensioni (30\,000 documenti): senza indice MongoDB esegue una scansione completa (\texttt{COLLSCAN}) su tutti i 30k documenti, con indice risolve la stessa query in $O(\log n)$ (\texttt{IXSCAN}). Le altre query mostrano guadagni modesti (1.1--1.2$\times$) perché le collezioni coinvolte sono piccole ($\le$ 356 documenti): a questa scala \texttt{COLLSCAN} e \texttt{IXSCAN} hanno tempi di esecuzione comparabili, e l'overhead della rete Docker domina la misura. Lo speedup di questi indici crescerebbe proporzionalmente al volume di dati.

\subsection{Scomposizione del costo della compilazione}
\label{sec:compile-bench}

Oltre al confronto con e senza indici, è stato strumentato l'endpoint di compilazione \texttt{POST /projects/\{id\}/compile} per stabilire \emph{dove} si concentri effettivamente il tempo di risposta. Il flusso è stato scomposto in tre fasi misurate separatamente con \texttt{time.perf\_counter()}:

\begin{enumerate}[leftmargin=2.2em]
  \item \textbf{Query al database} -- recupero dei sorgenti del progetto: \texttt{find(\{project\_id, file\_type: tex|bib\})} sulla collezione \texttt{files};
  \item \textbf{Scrittura su filesystem} -- copia dei sorgenti nella directory temporanea isolata;
  \item \textbf{Compilazione LaTeX} -- esecuzione del subprocess \texttt{tectonic}.
\end{enumerate}

Per garantire la riproducibilità della misura, il benchmark compila sempre lo stesso documento: un file \texttt{.tex} fisso (\texttt{seed/fixtures/benchmark.tex}) inserito dallo script di popolazione in un progetto dedicato. Un documento generato casualmente avrebbe infatti un costo di compilazione variabile, rendendo le medie non confrontabili tra esecuzioni. La metodologia è la stessa adottata per gli indici: media su più esecuzioni con una run di \emph{warm-up} esclusa, necessaria in questo caso perché Tectonic scarica e mette in cache i pacchetti alla prima invocazione.

\begin{table}[h!]
\centering
\small
\begin{tabular}{@{}lrr@{}}
\toprule
\textbf{Fase} & \textbf{Tempo medio (ms)} & \textbf{\% del totale} \\
\midrule
Query al database & 2.977 & 0.69\,\% \\
Scrittura su filesystem & 0.192 & 0.04\,\% \\
Compilazione LaTeX (Tectonic) & 428.834 & 99.27\,\% \\
\midrule
\textbf{Totale} & \textbf{432.004} & \textbf{100.00\,\%} \\
\bottomrule
\end{tabular}
\caption{Ripartizione del tempo di compilazione, media su 5 esecuzioni misurate (una run di warm-up esclusa), documento fisso.}
\end{table}

\textbf{Interpretazione.} Il livello dati incide per lo \textbf{0.69\,\%} del tempo totale di risposta: anche azzerando completamente il costo della query, il tempo complessivo passerebbe da 432 a 429 ms, un miglioramento non percepibile dall'utente. Si tratta di un'applicazione diretta della legge di Amdahl: lo speedup ottenibile ottimizzando una componente è limitato superiormente dalla frazione di tempo che tale componente occupa. Il collo di bottiglia è la compilazione vera e propria, che è lavoro di CPU e I/O esterno al database e non comprimibile agendo sul modello dei dati.

È comunque utile osservare che la query di compilazione (2.977 ms) risulta più costosa della corrispondente query sui file misurata in Tabella~10 (1.609 ms), pur insistendo sulla stessa collezione e sullo stesso indice. La differenza è nel tipo di accesso: la seconda è un \texttt{count\_documents} che si risolve interamente sull'indice, mentre la prima è una \texttt{find} che materializza i documenti completi, incluso il campo \texttt{content} contenente l'intero sorgente LaTeX. È la distinzione tra la lettura delle sole chiavi di indice e il successivo accesso ai documenti.

Le leve disponibili sul lato database sarebbero due: una \emph{projection} esplicita, per limitare i campi trasferiti a quelli effettivamente utilizzati, e un indice composto \texttt{project\_id + file\_type}, dato che la query filtra su entrambi i campi mentre l'indice attuale copre solo il primo. Entrambe ridurrebbero i 2.977 ms misurati, ma -- coerentemente con quanto sopra -- l'effetto sul tempo di risposta percepito resterebbe sotto la soglia di misurabilità. La scelta di non implementarle è quindi una conseguenza documentata della misura, non un'omissione.

% =====================================================
\section{Aggregation pipeline (JOIN)}
% =====================================================

L'endpoint \texttt{GET /stats/projects} realizza una pipeline di aggregazione multi-stadio che unisce quattro collezioni, rappresentando la principale dimostrazione di JOIN cross-collection in MongoDB.

\begin{figure}[h!]
\centering
\begin{tikzpicture}[
  stage/.style={draw, rounded corners, minimum width=4.4cm, minimum height=0.9cm, align=center, fill=orange!10},
  coll/.style={draw, rounded corners, minimum width=3cm, minimum height=0.8cm, align=center, fill=blue!8},
  arr/.style={-Stealth, thick}
]
\node[coll] (proj) {\texttt{projects}};
\node[stage, below=0.5cm of proj] (l1) {\$lookup $\to$ \texttt{users} (owner)};
\node[stage, below=0.5cm of l1] (l2) {\$lookup $\to$ \texttt{permissions}};
\node[stage, below=0.5cm of l2] (l3) {\$lookup $\to$ \texttt{activity\_logs}};
\node[stage, below=0.5cm of l3] (grp) {\$group by \texttt{department}};
\node[stage, below=0.5cm of grp] (srt) {\$sort by \texttt{total\_activity} DESC};

\draw[arr] (proj) -- (l1);
\draw[arr] (l1) -- (l2);
\draw[arr] (l2) -- (l3);
\draw[arr] (l3) -- (grp);
\draw[arr] (grp) -- (srt);
\end{tikzpicture}
\caption{Pipeline di aggregazione di \texttt{/stats/projects}. Filtri opzionali (\texttt{department}, \texttt{date\_from}, \texttt{date\_to}) sono applicati in uno stadio \$match prima dei lookup.}
\end{figure}

\subsection{Esempio di chiamata e risposta}

\begin{lstlisting}
GET /stats/projects?department=Engineering

[
  {
    "department": "Engineering",
    "project_count": 12,
    "total_collaborators": 34,
    "avg_collaborators": 2.83,
    "total_activity": 1520
  }
]
\end{lstlisting}

\pagebreak

% =====================================================
\section{Metriche di progetto}
% =====================================================

\begin{table}[h!]
\centering
\begin{tabular}{@{}lr@{}}
\toprule
\textbf{Area} & \textbf{Righe di codice (LOC)} \\
\midrule
Backend Python (\texttt{backend/}) & 1\,245 \\
Frontend JS/JSX (\texttt{frontend/src/}) & 1\,466 \\
Script seed + scripts (indici/benchmark/validazione) & 1\,024 \\
\midrule
\textbf{Totale} & \textbf{3\,735} \\
\bottomrule
\end{tabular}
\caption{Conteggio righe di codice (\texttt{wc -l} sui file \texttt{.py}/\texttt{.js}/\texttt{.jsx}, esclusi \texttt{node\_modules}).}
\end{table}

\begin{table}[h!]
\centering
\begin{tabular}{@{}lr@{}}
\toprule
\textbf{Metrica} & \textbf{Valore} \\
\midrule
Collezioni MongoDB & 5 \\
Indici & 6 (5 di query + 1 univoco) \\
Router API & 7 \\
Endpoint REST & 21 \\
Ruoli RBAC & 3 (Admin / Editor / Viewer) \\
Pagine frontend & 6 \\
Componenti UI condivisi & 4 \\
Issue chiuse & 14 \\
Pull request mergiate & 14 \\
% TODO: ore/persona spese sul progetto (da form di valutazione)
% TODO: numero di contributor attivi
\bottomrule
\end{tabular}
\caption{Metriche riassuntive di progetto.}
\end{table}

% =====================================================
\section{Uso dell'intelligenza artificiale}
% =====================================================

In conformità con i requisiti del corso, l'uso di strumenti di intelligenza artificiale generativa è dichiarato in modo trasparente. La tracciatura puntuale, issue per issue e PR per PR, è contenuta nei form di valutazione compilati durante lo svolgimento del progetto (FORM~3 e FORM~4). Di seguito una sintesi delle aree in cui gli strumenti sono stati impiegati.

\begin{itemize}[leftmargin=2.2em]
  \item \textbf{Progettazione del modello dati} (Issue \#1) -- supporto al ragionamento sui criteri di embedding e referencing e sulla gestione dei file multimediali, con esplicitazione dei vincoli (limite di 16\,MB per documento) alla base delle scelte adottate.
  \item \textbf{Generazione dati sintetici} (Issue \#2) -- supporto alla stesura dello script di seed; output verificato ed eseguito prima del merge.
  \item \textbf{Query e aggregazioni} (Issue \#7) -- supporto alla stesura della pipeline di aggregazione multi-stadio (\texttt{\$lookup}, \texttt{\$group}, \texttt{\$sort}); risultati verificati eseguendo la pipeline sui dati di seed.
  \item \textbf{Indici e benchmark} (Issue \#8--\#9) -- supporto alla definizione degli indici, alla lettura dell'output di \texttt{explain()} e alla struttura dello script di misura; metodologia e interpretazione dei risultati definite dal team.
  \item \textbf{Autenticazione e compilazione} (Issue \#20--\#21) -- supporto alla configurazione dell'hashing con \texttt{bcrypt} e dei token JWT, e all'invocazione isolata del compilatore esterno (directory temporanea, timeout, semaforo).
  \item \textbf{Strumentazione della compilazione} (Sezione~\ref{sec:compile-bench}) -- supporto all'implementazione della misurazione per fasi e del relativo endpoint di benchmark; metodologia (documento fisso, warm-up escluso) e interpretazione dei risultati definite dal team.
  \item \textbf{Documentazione} (Issue \#10) -- Claude Code usato per la generazione della documentazione, il fix del Dockerfile (build-arg \texttt{TARGETARCH}), le traduzioni e il polish del frontend.
\end{itemize}

In tutti i casi l'output generato è stato eseguito e verificato prima dell'integrazione, e revisionato da un membro del team diverso dall'autore, secondo il workflow standard: 1 issue $\to$ 1 branch $\to$ 1 PR $\to$ 1 review $\to$ merge. Nessun contenuto è stato integrato senza essere compreso e verificabile dal team.

% =====================================================
\section{Conclusioni}
\label{sec:conclusioni}
% =====================================================

Il progetto OpenTeX è stato sviluppato attraverso 14 issue, ciascuna su un branch dedicato e chiusa da una pull request collegata tramite \texttt{Closes}. Le prime dieci coprono i requisiti del progetto; le ultime quattro sono estensioni realizzate a partire dai \emph{nice to have} della proposta iniziale.

\begin{table}[h!]
\centering
\scriptsize
\begin{tabular}{@{}clp{7.6cm}c@{}}
\toprule
\textbf{Issue} & \textbf{Branch} & \textbf{Titolo} & \textbf{PR} \\
\midrule
\multicolumn{4}{@{}l}{\textit{Requisiti di progetto}} \\
\#1 & \texttt{issue-1-schema-validation} & Schema e validation delle 5 collezioni MongoDB & \#11 \\
\#2 & \texttt{issue-2-seed-data} & Generatore di dati sintetici (seed) & \#12 \\
\#3 & \texttt{issue-3-docker-setup} & Setup repository, Docker Compose e configurazione base & \#13 \\
\#4 & \texttt{issue-4-projects-crud} & API CRUD per la gestione dei progetti & \#14 \\
\#5 & \texttt{issue-5-permissions} & Sistema di permessi e gestione collaboratori & \#18 \\
\#6 & \texttt{issue-6-dashboard-ui} & UI Dashboard ``I miei Progetti'' & \#19 \\
\#7 & \texttt{issue-7-aggregation} & Aggregation pipeline per statistiche (JOIN tra collezioni) & \#15 \\
\#8 & \texttt{issue-8-indexing} & Creazione e ottimizzazione indici & \#16 \\
\#9 & \texttt{issue-9-benchmark} & Benchmark query con e senza indici & \#17 \\
\#10 & \texttt{issue-10-documentation} & Documentazione finale e proof of concept & \#28 \\
\midrule
\multicolumn{4}{@{}l}{\textit{Estensioni oltre i requisiti minimi}} \\
\#20 & \texttt{issue-20-auth} & Autenticazione con credenziali e token JWT & \#24 \\
\#21 & \texttt{issue-21-latex-compiler} & Endpoint di compilazione LaTeX lato backend & \#25 \\
\#22 & \texttt{issue-22-editor-ui} & Editor multi-pannello in stile Overleaf & \#26 \\
\#23 & \texttt{issue-23-pdf-preview} & Pannello di anteprima del PDF compilato & \#27 \\
\bottomrule
\end{tabular}
\caption{Mappatura completa issue $\to$ branch $\to$ pull request.}
\end{table}

Tutte le issue risultano chiuse e tutte e 14 le pull request mergiate su \texttt{main} (l'ultima, PR \#28, relativa alla documentazione finale); nessuna modifica è entrata nel branch principale al di fuori di una pull request revisionata. Il sistema è interamente riproducibile da zero seguendo le istruzioni del README, non contiene credenziali reali nel repository, ed è stato verificato end-to-end (registrazione, login, creazione progetto, upload/editing file, compilazione LaTeX, gestione permessi, consultazione statistiche).

\end{document}